% ============================================================================
\section{Discussion and Limitations}
\label{sec:discussion}

\subsection{What the result changes}
Three practical consequences follow. First, a feedback gain measured on a random
split should be re-measured after removing near-duplicate procedures and reply
templates; otherwise the number describes the data regime as much as the method.
The familiar-panel blind effect is $-0.0444$; on strict history it is $+0.0008$.
A team that budgets on the first number will overstate both the risk and the
opportunity. Second, tuning the legacy baseline matters. The published Laplace cell
ranks 23rd of 25 on the strict panel; an untuned lift is not a fair comparator, and
a tuned one still does not produce a confirmed gain on strict history. Third,
similarity to a single recorded reply is not procedural correctness. The expert and
the cosine metric diverge on the familiar arm, and the expert's preference is
itself sensitive to answer length. Evaluation protocols should state the evidence
set, keep the prior aligned with it, and add an expert check with length control.

The result also changes the interpretation of the earlier sign reversal. The
reversal is real for the legacy Laplace lift on the familiar split: the same
reranking helps under a resolution-informed judge and harms under a ticket-only
judge. But the blind harm is not robust to the data boundary, and the conditioned
gain is small on strict history. The accurate summary is narrower than ``the judge
decides the sign'': an uncalibrated popularity lift can help under one judge and
hurt under another on overlapping data, while a calibrated lift removes the harm
without a confirmed gain on new procedures.

\subsection{What we do not claim}
We do not claim a deployable improvement. The calibrated blind gain is not
significant and lies below its run-specific detection threshold; the conditioned
protocol is a resolution-informed upper bound for historical pairs, not a claim that
a new ticket's resolution is available. The feedback is simulated by an LLM judge,
not collected from users. The findings come from one organizational corpus and one
language; external validity is untested. The primary quality metric is similarity
to one recorded reply, not verified resolution success, and per-ticket agreement
between metrics and the expert is limited even when aggregate directions agree.
The alternative-retriever comparison is matched on the feedback configuration but
still uses one corpus and one generator. The secondary analyses (scope grid, formula
ordering, expected-value selection, negative results) are exploratory and are not
used for confirmation. We do not claim that calibration and control are generally
substitutable; we claim that in the tested configurations the gate adds little once
the prior is calibrated, and that it pays mainly where always-on is harmful.

\subsection{Limitations}
The corpus is a single enterprise in a single language. The expert study uses 105
changed pairs, 35 per arm; with 2--7 discordant pairs per arm it can detect only
large differences, and one expert is not a population sample. The second expert is
in progress; until the sample is complete, no inter-rater agreement is reported.
Procedural correctness is not ticket resolution, and the automatic metric compares
to one recorded reply. Feedback is simulated by an LLM judge, not collected from
users. Routing labels are the recorded resolving team and the reply-derived intent
class, so the results are label-conditional; a production system would need a
predictor or a human routing step. Hosted-model drift is bounded by a frozen
generation regime and hashed artifacts, not eliminated. BGE and BERTScore are not
recomputed for every history-only panel; the independent check is the expert
assessment.

\subsection{Threats to validity}
\textbf{Construct validity.} Cosine similarity to a single reference reply is a
proxy for usefulness, not resolution success; a ticket may admit several valid
replies. The expert assessment reduces this threat but does not eliminate it: the
expert judges procedural correctness from the ticket and the two replies, and the
rubric is sensitive to answer length. The feedback judge is resolution-informed on
historical pairs, so it measures procedural equivalence to a recorded resolution,
not the usefulness of a reply to a new user.

\textbf{Internal validity.} The data regimes share the corpus and the generation
regime, which strengthens within-study contrasts but limits causal isolation. The
strict panel changes the data boundary and the candidate pool together; the
ablation isolates the lift factors, not every possible confound. The split audit
shows that the earlier ``disjoint'' split was not disjoint, but it does not prove
that every remaining difference between panels is due to overlap. The matched
retriever comparison holds the feedback configuration fixed, but it still uses one
generator and one corpus.

\textbf{External validity.} The corpus is one enterprise in one language; absolute
effect sizes should not be transferred without re-validation. The saturation
diagnostic is cheap and portable, and we release it so that others can compute it
on their own feedback sources. The practical recipe in Section~\ref{sec:discussion}
is a set of hypotheses to test, not a general law.

\textbf{Statistical validity.} The MDE is run-specific and states the effect sizes
the study could reliably detect; it is not an equivalence boundary. The strict
conditioned effect and the strict protocol contrast are below their MDEs, so we
describe them as small and imprecise. The familiar blind harm is above its MDE and
survives Holm correction; the strict conditioned effect does not survive Holm. The
two-blind-run agreement at four decimals in the earlier evaluation is a coincidence
of aggregation: the runs differ per ticket (mean absolute difference 0.0671,
median 0.0337, 11.8\,\% identical deltas, $r=.799$), and both inherit the saturated
prior.

\subsection{Practical implications}
This article suggests a short checklist for teams that use feedback reranking.
(1)~\emph{Measure overlap first}: report the share of evaluation tickets that share
a full reply, a reply opening, or a title+subcategory signature with training.
(2)~\emph{Align the prior with the allowed evidence}: estimate the scope base rate
from exactly the feedback rows that are admissible, and state the estimator.
(3)~\emph{Tune the baseline}: compare against a tuned Laplace cell under a matched
selection budget, not only against the inherited parameters. (4)~\emph{Add a reply-similarity baseline that uses no LLM judge}: if it matches the LLM judge, the judge may be buying
overlap rather than procedure. (5)~\emph{Check a sample with experts}: use blinded
A/B order, per-reply correctness, explicit ties and both-incorrect options, and a
length control. (6)~\emph{Treat a gate as risk control}: it pays where always-on is
harmful, not as a general source of gain. These are hypotheses to test on other
corpora, not a guarantee.
